\chapter{The multiperiod binomial market}
\label{ch:binomial-multi}

\begin{question}
Over many periods the stock moves up or down each period. How do we price a claim and copy it, step by step through the tree?
\end{question}

\begin{whatfor}
\begin{itemize}
  \item Self-financing strategies over $T$ periods.
  \item Prices any claim $g(S_T)$ by working backwards through the tree, one period at a time.
  \item Gets the closed formula: a binomial average of the payoff under $p^*$, discounted; and shows the discounted stock is a martingale under $P^*$.
\end{itemize}
\skipline{fs:binomial-multi}
\end{whatfor}

\begin{prereq}
From \crefrange{ch:binomial-one}{ch:fundamental}: binomial market, trading strategy, self-financing, perfect hedge, $p^*=\frac{1+r-d}{u-d}$, EMM, fair price. From Part~1: martingale. From Part~0: the binomial distribution.
\end{prereq}

%=====================================================================
\section{Strategies over many periods}

\begin{setup}[The $T$-period binomial market]
Bond $B_t=(1+r)^t$; stock $S_t=uS_{t-1}$ or $dS_{t-1}$, $t=1,\dots,T$, with $d<1+r<u$. An outcome is a path $\omega=(v_1,\dots,v_T)$, $v_t\in\{u,d\}$, and $S_t=v_1\cdots v_tS_0$. A claim pays $X=g(S_T)$ at $T$.
\end{setup}

A strategy holds $\Delta_t$ shares and $b_t$ bonds during period $t$, chosen at time $t-1$ from $S_0,\dots,S_{t-1}$. Its value is $V_t=\Delta_tS_t+b_t(1+r)^t$, with $V_0=\Delta_1S_0+b_1$. It is self-financing (\cref{ch:binomial-one}) when rebalancing costs nothing:
\[
  \Delta_tS_t+b_t(1+r)^t=\Delta_{t+1}S_t+b_{t+1}(1+r)^t,\qquad t=1,\dots,T-1 .
\]
It copies $X$ if $V_T=X$ on every path.

%=====================================================================
\section{Backward induction}

\begin{theorem}[Pricing in the $T$-period tree]
\label{thm:multi}
Write $V_t[s]$ for the value at time $t$ when $S_t=s$. Then $V_T[s]=g(s)$ and, going backwards,
\[
  V_t[s]=\frac{1}{1+r}\Big[p^*V_{t+1}[us]+(1-p^*)V_{t+1}[ds]\Big]=\E_{P^*}\Big(\frac{V_{t+1}}{1+r}\Bigm|S_t=s\Big),
\]
with the copying holding over period $t+1$
\[
  \Delta_{t+1}[s]=\frac{V_{t+1}[us]-V_{t+1}[ds]}{us-ds},\qquad b_{t+1}[s]=\frac{uV_{t+1}[ds]-dV_{t+1}[us]}{(1+r)^{t+1}(u-d)} .
\]
In closed form,
\[
  \begin{aligned}
  V_t[s]&=\frac{1}{(1+r)^{T-t}}\sum_{j=0}^{T-t}\binom{T-t}{j}(p^*)^j(1-p^*)^{T-t-j}\,g\big(u^jd^{T-t-j}s\big)\\
  &=\E_{P^*}\Big(\frac{g(S_T)}{(1+r)^{T-t}}\Bigm|S_t=s\Big).
  \end{aligned}
\]
The fair price of $X$ is $V_0[S_0]$.
\end{theorem}
\begin{proof}
Each period is the one-period problem of \cref{ch:binomial-one}.
\begin{steps}
  \item \emph{Last period.} At time $T-1$ with $S_{T-1}=s$, copy $g$ over one period: $\Delta_Tus+b_T(1+r)^T=g(us)$ and $\Delta_Tds+b_T(1+r)^T=g(ds)$. Solving gives the formulas with $t+1=T$, and the cost is $V_{T-1}[s]=\frac{1}{1+r}[p^*g(us)+(1-p^*)g(ds)]$, as in \cref{ch:binomial-one} (one bond costs $(1+r)^{T-1}$ at time $T-1$ and is worth $(1+r)^T$ at $T$).
  \item \emph{One step further back.} $V_{T-1}[\cdot]$ is now a known function of $S_{T-1}$: treat it as a new claim with expiry $T-1$ and repeat. The value received at time $T-1$ is exactly the cost of the next holding, so the strategy is self-financing.
  \item \emph{Closed form.} After $T-t$ backward steps, each path with $j$ ups and $T-t-j$ downs ends at $u^jd^{T-t-j}s$ and carries weight $(p^*)^j(1-p^*)^{T-t-j}$; there are $\binom{T-t}{j}$ such paths. \qedhere
\end{steps}
\end{proof}

\begin{example}[A two-period call]
\label{ex:two-period}
$S_0=100$, $u=1.2$, $d=0.8$, $r=0.05$, $T=2$, call with $K=100$; $p^*=\frac{1.05-0.8}{0.4}=0.625$.
\begin{steps}
  \item Time 2: $S_2=144,96,64$; payoffs $44,0,0$.
  \item Time 1: $V_1[120]=\frac{0.625\cdot44}{1.05}\approx26.19$, $V_1[80]=0$.
  \item Time 0: $V_0=\frac{0.625\cdot26.19}{1.05}\approx15.59$. Closed form: $\frac{0.625^2\cdot44}{1.05^2}=15.59$.
  \item Hedge at time 0: $\Delta_1=\frac{26.19-0}{120-80}\approx0.655$ shares and $V_0-\Delta_1S_0\approx-49.89$ in the bond. After an up-move, rebalance to $\Delta_2=\frac{44-0}{144-96}\approx0.917$ shares and $26.19-0.917\cdot120\approx-83.81$ in the bond. Self-financing check: before rebalancing, $0.655\cdot120-49.89\cdot1.05=26.19$.
\end{steps}
\end{example}

\begin{pic}
\begin{tikzpicture}[x=3.4cm,y=1.25cm,every node/.style={font=\small}]
  \node[draw,rounded corners,align=center] (a) at (0,0) {$S=100$\\ $V=15.59$};
  \node[draw,rounded corners,align=center,fill=defc!8] (b) at (1,1) {$S=120$\\ $V=26.19$};
  \node[draw,rounded corners,align=center,fill=trapc!6] (c) at (1,-1) {$S=80$\\ $V=0$};
  \node[draw,rounded corners,align=center] (d) at (2,2) {$S=144$\\ $V=44$};
  \node[draw,rounded corners,align=center] (e) at (2,0) {$S=96$\\ $V=0$};
  \node[draw,rounded corners,align=center] (f) at (2,-2) {$S=64$\\ $V=0$};
  \draw[-{Latex}] (a)--node[below right,font=\scriptsize]{$p^*=0.625$}(b); \draw[-{Latex}] (a)--node[below,font=\scriptsize]{$0.375$}(c);
  \draw[-{Latex}] (b)--(d); \draw[-{Latex}] (b)--(e); \draw[-{Latex}] (c)--(e); \draw[-{Latex}] (c)--(f);
  \node[thmc,font=\scriptsize] at (0,0.75) {$\Delta_1=0.655$};
  \node[thmc,font=\scriptsize] at (1,1.75) {$\Delta_2=0.917$};
  \draw[-{Latex},prec,thick] (2.45,-2.6)--(-0.2,-2.6) node[midway,below,font=\scriptsize] {values are filled in backwards, from the payoffs};
\end{tikzpicture}
\picnote{The two-period tree of \cref{ex:two-period}. Each node's value is the discounted $p^*$-average of the two values to its right; each node's $\Delta$ is the change in value divided by the change in stock price.}
\end{pic}

%=====================================================================
\section{The discounted stock is a martingale}

\begin{result}[Martingales under $P^*$]
\label{res:multi-mg}
Under $P^*$ the moves are independent, up with probability $p^*$. The discounted stock $S_t^*=(1+r)^{-t}S_t$ and the discounted value $(1+r)^{-t}V_t$ of any self-financing strategy are martingales; so $V_0=\E_{P^*}\big((1+r)^{-T}X\big)$ for any copied $X$.
\end{result}
\begin{proof}
$\E_{P^*}(S_{t+1}\given\F_t)=S_t(p^*u+(1-p^*)d)=(1+r)S_t$, so $\E_{P^*}(S^*_{t+1}\given\F_t)=S^*_t$ (the one-step test of \cref{ch:martingales}). The value: $V_{t+1}-V_t=\Delta_{t+1}(S_{t+1}-S_t)+b_{t+1}((1+r)^{t+1}-(1+r)^t)$ by self-financing, and discounting turns this into $\Delta_{t+1}(S^*_{t+1}-S^*_t)$, a stake chosen from the past on a martingale (\cref{ch:stopping}).
\end{proof}

\paragraph{The price as an exponential of a random walk.} Write $S_T=S_0e^{Z_T}$ with $Z_T=Y_1+\dots+Y_T$, where $Y_k=\ln u$ (prob.\ $p^*$) or $\ln d$ (prob.\ $1-p^*$) independently. Then
\[
  V_0=\E_{P^*}\Big(\frac{g(S_0e^{Z_T})}{(1+r)^T}\Big).
\]
Chapter~22 lets the number of periods grow with $u,d$ shrinking, and $Z_T$ becomes normal.

%=====================================================================
\formulasheet{fs:binomial-multi}

\begin{mastercard}[label=card:tree]{Pricing and hedging in the $T$-period binomial tree}
\cref{card:finite}, one period at a time, from the payoff backwards. $S_t=uS_{t-1}$ or $dS_{t-1}$, $p^*=\frac{1+r-d}{u-d}$, a claim pays $g(S_T)$; $V_t[s]$ is the value at time $t$ when $S_t=s$:
\[
  \boxed{\;V_t[s]=\frac{p^*\,V_{t+1}[us]+(1-p^*)\,V_{t+1}[ds]}{1+r},\qquad V_T[s]=g(s).\;}
\]
\begin{plugrows}
  \plugrow{Price today}{\pstep{Box: $V_t[s]=\frac{p^*V_{t+1}[us]+(1-p^*)V_{t+1}[ds]}{1+r}$, $V_T[s]=g(s)$.}\pstep{Start at $T$ with the payoffs; step back one period at a time to $t=0$.}\pstep{All $T$ steps at once: $j$ up-moves happen in $\binom Tj$ orders, each with weight $(p^*)^j(1-p^*)^{T-j}$.}}{$V_0=\dfrac{1}{(1+r)^T}\sum_{j=0}^{T}\binom Tj(p^*)^j(1-p^*)^{T-j}g\big(u^jd^{T-j}S_0\big)$}{$d<1+r<u$; European claims $g(S_T)$}{$S_0=100$, $u=1.2$, $d=0.8$, $r=0.05$, $T=2$, $K=100$: $p^*=\frac{1.05-0.8}{0.4}=0.625$; payoffs $44,0,0$ at $144,96,64$; $V_0=\frac{0.625^2(44)}{1.05^2}=\frac{17.19}{1.1025}=15.59$}
  \plugrow{One period}{\pstep{Box: $V_t[s]=\frac{p^*V_{t+1}[us]+(1-p^*)V_{t+1}[ds]}{1+r}$, with $T=1$, $t=0$, $V_1=g$.}\pstep{This is \cref{card:finite}, row 1: $X^*=\frac{p^*X_u+(1-p^*)X_d}{1+r}$.}}{the one-step price of \cref{card:finite}, row 1}{one period}{$S_0=1$, $u=1.75$, $d=0.5$, $r=0.25$, $K=1$: $p^*=0.6$; $\frac{0.6(0.75)+0.4(0)}{1.25}=0.36$}
  \plugrow{A call: only the ``in the money'' terms}{\pstep{Row 1 with $g(s)=(s-K)^+$: only $j\ge m$ pay, $m$ = least $j$ with $u^jd^{T-j}S_0>K$.}\pstep{Split $(S_T-K)$ into the $S_T$ part and the $K$ part.}\pstep{$S_T$ part: $\frac{(p^*u)^j((1-p^*)d)^{T-j}S_0}{(1+r)^T}$, which is $S_0\hat p^j(1-\hat p)^{T-j}$ with $\hat p=\frac{up^*}{1+r}$.}\pstep{$K$ part: $\frac K{(1+r)^T}\times$ the binomial probability.}}{$C=S_0\sum_{j\ge m}\binom Tj\hat p^{\,j}(1-\hat p)^{T-j}-\dfrac{K}{(1+r)^T}\sum_{j\ge m}\binom Tj(p^*)^j(1-p^*)^{T-j}$}{call, $g$ increasing; $\hat p$ is again a probability}{same: $m=2$, $\hat p=\frac{1.2(0.625)}{1.05}=0.714$: $100(0.714^2)-\frac{100}{1.1025}(0.625^2)=100(0.510)-90.70(0.391)=51.02-35.43=15.59$}
  \plugrow{A put}{\pstep{\cref{card:parity}, row 1, over $T$ periods: $P=C-S_0+KD_T$, $D_T=(1+r)^{-T}$.}\pstep{Take $C$ from row 1 or 3.}}{$P=C-S_0+K(1+r)^{-T}$}{European}{same: $\frac{100}{1.1025}=90.70$; $P=15.59-100+90.70=6.29$}
  \plugrow{The hedge at a node}{\pstep{At node $(t,s)$ the next period is one period of \cref{card:finite}, row 1.}\pstep{Copy: $\Delta=\frac{X_u-X_d}{(u-d)s}$ with $X_u=V_{t+1}[us]$, $X_d=V_{t+1}[ds]$.}\pstep{Bond money: what is left of the value, $V_t[s]-\Delta s$.}}{$\Delta=\dfrac{V_{t+1}[us]-V_{t+1}[ds]}{us-ds}$ shares and $V_t[s]-\Delta s$ in the bond}{use values one step ahead}{$V_1[120]=\frac{0.625(44)}{1.05}=26.19$, $V_1[80]=0$; time $0$: $\Delta_1=\frac{26.19-0}{120-80}=0.655$, bond $15.59-65.48=-49.89$; at $120$: $\Delta_2=\frac{44-0}{144-96}=0.917$, bond $26.19-110=-83.81$}
  \plugrow{Is the strategy self-financing?}{\pstep{\cref{card:selffin}, row 2: old holding at the new prices $=$ new holding at the new prices.}\pstep{Old holding at time 1: $\Delta_1S_1+(\text{bond})(1+r)$.}\pstep{New holding: $V_1[S_1]$ by row 5.}}{yes for the backward-induction hedge}{compare at the same price}{at $S_1=120$: $0.655(120)-49.89(1.05)=78.57-52.38=26.19=V_1[120]$}
  \plugrow{A discounted martingale}{\pstep{Box applied to the stock itself: $S_t=\frac{p^*uS_t+(1-p^*)dS_t}{1+r}$, i.e.\ $\E_{P^*}(S_{t+1}\given\F_t)=(1+r)S_t$.}\pstep{Divide by $(1+r)^{t+1}$: $\E_{P^*}\big((1+r)^{-(t+1)}S_{t+1}\given\F_t\big)=(1+r)^{-t}S_t$.}\pstep{The box says the same for $V_t$.}}{$(1+r)^{-t}S_t$ and $(1+r)^{-t}V_t$ are $P^*$-martingales}{self-financing strategies}{$0.625(120)+0.375(80)=75+30=105=1.05(100)$}
\end{plugrows}
\end{mastercard}
